%% kapa: exam=1 round=25 year=2014 period=1 form=A subject=economics subject_ko=경제원론 seq=3 number=43
%% source: Q-Net 감정평가사 기출문제 (cst00309) artlSeq=5101207 file=extracted/1008/5101207/2014년 제25회 감정평가사 1차시험 기출문제/1교시(A형).hwp
%% conversion: kordoc 4.24.51 -> tools/gen1.py (auto); review=proofread-sonnet; flags=none
\begin{question}{43}
효용을 극대화하는 소비자 甲이 소득 100으로 완전보완재인 $X$재와 $Y$재만을 소비하고 있다. $Y$재의 가격은 10으로 일정하고 $X$재의 가격이 40에서 10으로 하락할 때 소비자의 효용이 증가한다. 이 때 $X$재의 가격이 하락하는 대신 소득이 얼마나 증가해야 동일한 효용의 증가를 가져오는가를 나타내는 동등변이(equivalent variation)는 얼마인가?
\begin{choices}
  \item 50
  \item 100
  \item 150
  \item 200
  \item 250
\end{choices}
\end{question}
\answerall
\answernote{정답표에 1,2,3,4,5 모두 정답으로 표기(전원 정답)}
